\documentclass[10pt,a4paper]{amsart}
\usepackage{amsmath,amssymb,amsthm,graphicx}
% Retained preamble for the Stacks Project chapter morphisms.tex (Morphisms of Schemes).
% The chapter file pinned-original-source.tex carries only \input{preamble} on its first line, so no
% preamble text lives inside the chapter itself.  The Stacks Project uses one single shared preamble for
% every chapter; the snapshot stages it verbatim in sources/stacks-chow/preamble.tex.  The block below is
% copied mechanically, line for line, from that file's lines 14-18 (the Xy-pic setup) and lines 154-217
% (the native theorem declarations together with the mathematical macros).  Nothing was added, removed or
% redefined.  The wrapper class and the amsmath/amssymb/amsthm/graphicx base packages are supplied by the
% assembler, so the class/\externaldocument part of the shared preamble is not repeated here.
% For commutative diagrams we use Xy-pic
\usepackage[all]{xy}

% We use 2cell for 2-commutative diagrams.
\xyoption{2cell}
% Theorem environments.
%
\theoremstyle{plain}
\newtheorem{theorem}[subsection]{Theorem}
\newtheorem{proposition}[subsection]{Proposition}
\newtheorem{lemma}[subsection]{Lemma}

\theoremstyle{definition}
\newtheorem{definition}[subsection]{Definition}
\newtheorem{example}[subsection]{Example}
\newtheorem{exercise}[subsection]{Exercise}
\newtheorem{situation}[subsection]{Situation}

\theoremstyle{remark}
\newtheorem{remark}[subsection]{Remark}
\newtheorem{remarks}[subsection]{Remarks}

\numberwithin{equation}{subsection}

% Macros
%
\def\lim{\mathop{\mathrm{lim}}\nolimits}
\def\colim{\mathop{\mathrm{colim}}\nolimits}
\def\Spec{\mathop{\mathrm{Spec}}}
\def\Hom{\mathop{\mathrm{Hom}}\nolimits}
\def\Ext{\mathop{\mathrm{Ext}}\nolimits}
\def\SheafHom{\mathop{\mathcal{H}\!\mathit{om}}\nolimits}
\def\SheafExt{\mathop{\mathcal{E}\!\mathit{xt}}\nolimits}
\def\Sch{\mathit{Sch}}
\def\Mor{\mathop{\mathrm{Mor}}\nolimits}
\def\Ob{\mathop{\mathrm{Ob}}\nolimits}
\def\Sh{\mathop{\mathit{Sh}}\nolimits}
\def\NL{\mathop{N\!L}\nolimits}
\def\CH{\mathop{\mathrm{CH}}\nolimits}
\def\proetale{{pro\text{-}\acute{e}tale}}
\def\etale{{\acute{e}tale}}
\def\QCoh{\mathit{QCoh}}
\def\Ker{\mathop{\mathrm{Ker}}}
\def\Im{\mathop{\mathrm{Im}}}
\def\Coker{\mathop{\mathrm{Coker}}}
\def\Coim{\mathop{\mathrm{Coim}}}

% Boxtimes
%
\DeclareMathSymbol{\boxtimes}{\mathbin}{AMSa}{"02}

%
% Macros for moduli stacks/spaces
%
\def\QCohstack{\mathcal{QC}\!\mathit{oh}}
\def\Cohstack{\mathcal{C}\!\mathit{oh}}
\def\Spacesstack{\mathcal{S}\!\mathit{paces}}
\def\Quotfunctor{\mathrm{Quot}}
\def\Hilbfunctor{\mathrm{Hilb}}
\def\Curvesstack{\mathcal{C}\!\mathit{urves}}
\def\Polarizedstack{\mathcal{P}\!\mathit{olarized}}
\def\Complexesstack{\mathcal{C}\!\mathit{omplexes}}
% \Pic is the operator that assigns to X its picard group, usage \Pic(X)
% \Picardstack_{X/B} denotes the Picard stack of X over B
% \Picardfunctor_{X/B} denotes the Picard functor of X over B
\def\Pic{\mathop{\mathrm{Pic}}\nolimits}
\def\Picardstack{\mathcal{P}\!\mathit{ic}}
\def\Picardfunctor{\mathrm{Pic}}
\def\Deformationcategory{\mathcal{D}\!\mathit{ef}}

% --- appended 2026-10-12 (Chief): the project own definitions of the reference/slogan/history
% display environments, copied mechanically from the shared Stacks preamble lines 8-12
% (\usepackage{verbatim} plus the three \newenvironment lines, which the chapter preamble keeps
% verbatim from the project preamble). Nothing added, removed or redefined: these environments
% carry no mathematical content and the project itself typesets them as comments. Without these
% three lines a counted statement whose own text uses them fails with Environment slogan undefined.
\usepackage{verbatim}
\newenvironment{reference}{\comment}{\endcomment}
\newenvironment{slogan}{\comment}{\endcomment}
\newenvironment{history}{\comment}{\endcomment}
\usepackage{cleveref}
\crefname{originallemma}{original lemma}{original lemmas}
\Crefname{originallemma}{Original lemma}{Original lemmas}
\makeatletter
\newlabel{lemma-image-nowhere-dense-finite}{{03HX}{source}{}{sourcepointer-1}{}}
\newlabel{lemma-image-nowhere-dense-finite@cref}{{[originallemma][0][]03HX}{[1][0][]source}{}{}{}}
\makeatother
\crefname{originallemma}{original lemma}{original lemmas}
\Crefname{originallemma}{Original lemma}{Original lemmas}
\makeatletter
\newlabel{lemma-quasi-finite-affine}{{03GU}{source}{}{sourcepointer-2}{}}
\newlabel{lemma-quasi-finite-affine@cref}{{[originallemma][0][]03GU}{[1][0][]source}{}{}{}}
\makeatother
\crefname{originallemma}{original lemma}{original lemmas}
\Crefname{originallemma}{Original lemma}{Original lemmas}
\makeatletter
\newlabel{topology-lemma-image-nowhere-dense-open}{{03J0}{source}{}{sourcepointer-3}{}}
\newlabel{topology-lemma-image-nowhere-dense-open@cref}{{[originallemma][0][]03J0}{[1][0][]source}{}{}{}}
\makeatother
\crefname{originallemma}{original lemma}{original lemmas}
\Crefname{originallemma}{Original lemma}{Original lemmas}
\makeatletter
\newlabel{topology-lemma-nowhere-dense}{{03HO}{source}{}{sourcepointer-4}{}}
\newlabel{topology-lemma-nowhere-dense@cref}{{[originallemma][0][]03HO}{[1][0][]source}{}{}{}}
\makeatother
\title{The image of a closed nowhere dense subset under a quasi-finite morphism is nowhere dense}
\author{The Stacks Project Authors}
\date{stacks-project:morphisms.tex:local-snapshot-20261012; GNU FDL 1.2 or later}
\begin{document}
\maketitle
\noindent\textit{Editorial note.} All delivered mathematical text is the author's own, copied from the pinned Stacks Project chapter file pinned-original-source.tex; no statement or proof step was rewritten, reordered or removed.\par\medskip
\noindent\textit{Editorial note.} Following the multi-claim convention for this source, source\_url is the source-level URL https://stacks.math.columbia.edu/ and the source label lemma-image-nowhere-dense-quasi-finite is carried in claim\_key and here rather than in a per-claim tag URL.\par\medskip
\noindent\textit{Editorial note.} The preamble is the chapter's shared Stacks Project preamble, retained as sources/stacks-morphisms/retained-preamble.tex: a line-for-line excerpt of the shared Stacks preamble staged with this snapshot. No macro was invented or redefined, and this card adds no preamble text of its own.\par\medskip
\noindent\textit{Editorial note.} Statement span 15253-15258 and proof span 15260-15284 come from the halfA pool card operations/pool-cards/stacks-morphisms-halfA.json and were re-read in the pinned chapter file before this card was written.\par\medskip
\noindent\textit{Editorial note.} The statement contains no hypotheses on the base beyond the morphism being quasi-finite; the author's formulation is preserved.\par\medskip
\noindent\textit{Editorial note.} The proof contains an xymatrix triangle displaying the factorization; it is copied verbatim and the retained chapter preamble loads Xy-pic.\par\medskip
\noindent\textit{Editorial note.} 4 cross-references occur inside the counted spans; they are listed in external\_references as navigation pointers only.\par\medskip
\noindent\textit{Allowed standard prerequisites (editorial).}\begin{itemize}
\item The same statement for finite morphisms, established earlier in the same chapter (and by a previous card of this lane). Cited at proof lines 15260-15261 and 15283-15284. Gives the conclusion in the finite case, to which the proof reduces.
\item The affine structure theorem for quasi-finite morphisms, factoring the morphism as an open immersion followed by a finite morphism, cited from the same chapter. Cited at proof lines 15270-15274. Produces the affine intermediate scheme Z over which the morphism is finite.
\item Nowhere density is local on the target, a finite union of closed nowhere dense subsets is nowhere dense, and the closure of the image of a closed nowhere dense subset under an open immersion is nowhere dense. Cited at proof lines 15262-15269 and 15276-15278. Allows the reduction to an affine base, the reduction to one member of the affine cover, and the passage from the subset to its closure in the intermediate scheme.
\end{itemize}\medskip
\ifdefined\hypertarget\hypertarget{sourcepointer-1}{}\fi
\noindent\textit{Source reference.} Original lemma 03HX: focused-seeds/AGP/sources/stacks-snapshot-20261011/tags.txt, source lines 4583--4583. Same-chapter reference used in the retained author text at source line(s) 15261, 15283. It is supplied only as a navigation pointer so that the retained text compiles; the counted item delivers the author's own statement and the author's complete proof, and no part of the referenced item is copied in or substituted for the counted argument.\par\medskip
\ifdefined\hypertarget\hypertarget{sourcepointer-2}{}\fi
\noindent\textit{Source reference.} Original lemma 03GU: focused-seeds/AGP/sources/stacks-snapshot-20261011/tags.txt, source lines 4541--4541. Same-chapter reference used in the retained author text at source line(s) 15270, 15522. It is supplied only as a navigation pointer so that the retained text compiles; the counted item delivers the author's own statement and the author's complete proof, and no part of the referenced item is copied in or substituted for the counted argument.\par\medskip
\ifdefined\hypertarget\hypertarget{sourcepointer-3}{}\fi
\noindent\textit{Source reference.} Original lemma 03J0: focused-seeds/AGP/sources/stacks-snapshot-20261011/tags.txt, source lines 4622--4622. Same-chapter reference used in the retained author text at source line(s) 15279. It is supplied only as a navigation pointer so that the retained text compiles; the counted item delivers the author's own statement and the author's complete proof, and no part of the referenced item is copied in or substituted for the counted argument.\par\medskip
\ifdefined\hypertarget\hypertarget{sourcepointer-4}{}\fi
\noindent\textit{Source reference.} Original lemma 03HO: focused-seeds/AGP/sources/stacks-snapshot-20261011/tags.txt, source lines 4574--4574. Same-chapter reference used in the retained author text at source line(s) 12547, 15267. It is supplied only as a navigation pointer so that the retained text compiles; the counted item delivers the author's own statement and the author's complete proof, and no part of the referenced item is copied in or substituted for the counted argument.\par\medskip
\begin{lemma}
\label{lemma-image-nowhere-dense-quasi-finite}
Let $f : Y \to X$ be a quasi-finite morphism of schemes.
Let $T \subset Y$ be a closed nowhere dense subset of $Y$.
Then $f(T) \subset X$ is a nowhere dense subset of $X$.
\end{lemma}

\begin{proof}
As in the proof of Lemma \ref{lemma-image-nowhere-dense-finite} this
reduces immediately to the case where the base $X$ is affine.
In this case $Y = \bigcup_{i = 1, \ldots, n} Y_i$ is a finite union
of affine opens (as $f$ is quasi-compact). Since each $T \cap Y_i$
is nowhere dense, and since a finite union of nowhere dense sets is
nowhere dense (see
Topology, Lemma \ref{topology-lemma-nowhere-dense}),
it suffices to prove that the image $f(T \cap Y_i)$ is nowhere dense in $X$.
This reduces us to the case where both $X$ and $Y$ are affine. At this point
we apply Lemma \ref{lemma-quasi-finite-affine} above to get a diagram
$$
\xymatrix{
Y \ar[rd]_f \ar[rr]_j & & Z \ar[ld]^\pi \\
& X &
}
$$
with $Z$ affine, $\pi$ finite and $j$ an open immersion.
Set $\overline{T} = \overline{j(T)} \subset Z$. By
Topology, Lemma \ref{topology-lemma-image-nowhere-dense-open}
we see $\overline{T}$ is nowhere dense in $Z$.
Since $f(T) \subset \pi(\overline{T})$
the lemma follows from the corresponding result in the finite case, see
Lemma \ref{lemma-image-nowhere-dense-finite}.
\end{proof}

\end{document}
